\appendix
\section{Additional Descriptive Results}

\begin{table}[!ht]
\caption{Recorded hardware and model settings. GPU and host peaks are observed
over the matched held-out stage; one-second sampling can miss short peaks.}
\label{tab:hardware}
\centering\small
\begin{tabular}{ll}
\toprule
Component & Recorded setting \\
\midrule
GPU & RTX 5090; 31.36 GiB \\
Primary model & Qwen3-8B; BF16; CUDA only \\
Model revision & \texttt{b968826d9c46} \\
Inference & vLLM 0.10.2; PyTorch 2.8.0+cu128 \\
Database & Neo4j 5.26.0; SQLite 3.45.1 WAL \\
Input / output limit & 4,096 / 1,536 tokens \\
Batch concurrency / interactive & 4 / 1 \\
CPU quota / memory limit & 27.2 CPU equivalents / 94.0 GB \\
Test GPU memory peak & 28,441.0 MiB \\
Test GPU utilization mean & 88.7\% \\
Test host cgroup memory peak & 28.5 GiB \\
\bottomrule
\end{tabular}

\end{table}

\begin{table}[!ht]
\caption{Generation behavior and supported-reference rates. Failed and truncated
columns are counts out of 240 cases per row; repairs count additional requests.
Reference validity requires support under the strict query contract.}
\label{tab:failures}
\centering\small\setlength{\tabcolsep}{3.5pt}
\begin{tabular}{llrrrrr}
\toprule
Source & Method & Failed & Truncated & Repairs & Abstain\% & Valid refs\% \\
\midrule
Synthetic & B0 & 1 & 0 & 0 & 0.4 & 48.9 \\
Synthetic & B1 & 0 & 0 & 127 & 8.3 & 100.0 \\
Synthetic & B2 & 0 & 0 & 127 & 8.3 & 100.0 \\
Synthetic & M1 & 0 & 0 & 127 & 8.3 & 100.0 \\
Synthetic & B3 & 0 & 0 & 127 & 8.3 & 100.0 \\
WESAD & B0 & 1 & 0 & 0 & 0.4 & 52.2 \\
WESAD & B1 & 0 & 0 & 141 & 8.3 & 100.0 \\
WESAD & B2 & 0 & 0 & 140 & 8.3 & 100.0 \\
WESAD & M1 & 0 & 0 & 141 & 8.3 & 100.0 \\
WESAD & B3 & 0 & 0 & 142 & 8.3 & 100.0 \\
PPG-DaLiA & B0 & 1 & 0 & 0 & 0.4 & 52.9 \\
PPG-DaLiA & B1 & 0 & 0 & 137 & 8.3 & 100.0 \\
PPG-DaLiA & B2 & 0 & 0 & 136 & 8.3 & 100.0 \\
PPG-DaLiA & M1 & 0 & 0 & 138 & 8.3 & 100.0 \\
PPG-DaLiA & B3 & 0 & 0 & 137 & 8.3 & 100.0 \\
\bottomrule
\end{tabular}

\end{table}

\begin{table}[!ht]
\caption{Post-hoc contract-scope breakdown. The last column counts direct numerical
observations with only interval-related errors that pass the existing exact checks
at their own stated interval and original knowledge time. These source-faithful
background facts remain errors in the unchanged primary metric. This narrow
breakdown does not exhaust the consequences of interval scope, including failures
propagated to dependent comparisons.}
\label{tab:contractscope}
\centering\small\setlength{\tabcolsep}{3pt}
\begin{tabular}{lrr}
\toprule
Source & B0 contract-invalid claims & Locally supported off-target facts \\
\midrule
Synthetic & 254 & 107 \\
WESAD & 278 & 115 \\
PPG-DaLiA & 267 & 111 \\
\bottomrule
\end{tabular}

\end{table}

\begin{table}[!ht]
\caption{Separate graph-versus-relational comparison of generated outcomes.
Differences are M1--B3 in percentage points with pointwise 95\% paired intervals;
N (S/E) gives retained subjects/episodes.}
\label{tab:relational}
\centering\small
\begin{tabular}{llrrl}
\toprule
Source & Outcome & N (S/E) & M1--B3 (pp) & 95\% interval (pp) \\
\midrule
Synthetic & Case error & 10/20 & 0.00 & [0.00, 0.00] \\
Synthetic & Fact recall & 10/20 & 0.00 & [0.00, 0.00] \\
Synthetic & Correction & 10/20 & 0.00 & [0.00, 0.00] \\
WESAD & Case error & 10/20 & 0.00 & [0.00, 0.00] \\
WESAD & Fact recall & 10/20 & 0.25 & [0.00, 0.75] \\
WESAD & Correction & 10/20 & 0.00 & [0.00, 0.00] \\
PPG-DaLiA & Case error & 10/20 & 0.00 & [0.00, 0.00] \\
PPG-DaLiA & Fact recall & 10/20 & -0.25 & [-0.75, 0.00] \\
PPG-DaLiA & Correction & 10/20 & 0.00 & [0.00, 0.00] \\
\bottomrule
\end{tabular}

\end{table}

\begin{table}[!ht]
\caption{Descriptive completion-time adequacy. Each explanation is checked against
its request snapshot and the latest ingested event when it completes. ``Became
stale'' counts request-adequate answers that no longer satisfy the two required
numerical slots. These counts do not replace the frozen metrics.}
\label{tab:staleness}
\centering\small\setlength{\tabcolsep}{2.4pt}
\begin{tabular}{rlrrrr}
\toprule
Events/s & Method & Calls & Input adequate & Completion adequate & Became stale \\
\midrule
1 & M1 & 5 & 0 & 0 & 0 \\
1 & B3 & 5 & 0 & 0 & 0 \\
5 & B3 & 23 & 0 & 0 & 0 \\
5 & M1 & 23 & 0 & 0 & 0 \\
20 & M1 & 90 & 0 & 0 & 0 \\
20 & B3 & 90 & 0 & 0 & 0 \\
\bottomrule
\end{tabular}

\end{table}
